\section{Smooth truncation of clock histories}\label{sec:history-cutoffs}

The replacement of the coupling variables requires bounds for clock paths
whose duration is of order $T=n^{2/3}$.  We obtain them by capping transition
densities on every scale of a binary time tree.  This section constructs the
caps, proves that they discard negligible Gaussian mass, and records the
conditional path laws to which the score estimates will apply.

All times in Sections~\ref{sec:history-cutoffs}--\ref{sec:clock-scores} are
in site-clock units.  A stationary query at lag $t$ is one path of length
$tT$ started from $\mu$.  A quench query $(s,t)$ is one path started from
$\nu_n$, with consecutive legs of lengths $sT$ and $tT$.  Fix a compact set
of query times.  Every complete leg then has length in $[cT,CT]$, where
$0<c<C<\infty$ are fixed.  We use Taylor orders at most
$K=24$.  The constant $\beta>0$ will be chosen sufficiently small after
$K$.  The notation $a\lesssim_\beta b$ means
$a\le C_0 n^{C_1\beta}b$, with $C_1$ depending at most on $K$ and the fixed
query set, but not on $\beta$.  Fixed logarithmic factors can be included in
this convention.  Moment orders and the precision of probability estimates
may be increased without changing $C_1$; their constants and the lower bound
on $n$ may change.  For a vector indexed by
$\mathscr E_n=\{\{i,j\}:1\le i<j\le n\}$, put
\[
 |v|_p=\left(\frac1{\binom n2}\sum_{e\in\mathscr E_n}|v_e|^p\right)^{1/p},
 \qquad |v|_\infty=\max_{e\in\mathscr E_n}|v_e|.
\]

\subsection{A static support valid for every interaction matrix}

Write $\theta_{ij}=J_{ij}/\sqrt n$.  In entry replacement, an order-$k$
derivative in $\theta_e$ is multiplied by $n^{-k/2}$; summing over the
order-$n^2$ entries gives the factor $n^{2-k/2}$.  The Gaussian and
Rademacher laws match through order three.  The averaged fourth derivative
of the truncated trace tests must therefore tend to zero, while higher
derivatives may grow within the remaining negative power of $n$.  The
estimates below are organized around this fourth-order requirement;
$K=24$ supplies a global remainder bound.

The Gibbs probability $\mu$ is strictly positive for every finite matrix
$\theta$.  For independent spins
$X^1,\ldots,X^4$ define
\[
 \rho_2(X^a,X^b)=\frac1n\sum_iX_i^aX_i^b,
 \qquad
 \rho_4(X^1,X^2,X^3,X^4)=\frac1n\sum_iX_i^1X_i^2X_i^3X_i^4.
\]
Let $\mathscr D_n$ be the dyadic numbers from $1$ through a fixed multiple
of $CT$, including the first dyadic number beyond that multiple.  Set
\begin{equation}\label{hist:static-scales}
 r_2(d)=n^{2\beta}d^{-1/2},\quad r_4(d)=n^{10\beta}d^{-3/4},\quad
 a_2(d)=n^{3\beta}nd^{-3/2},\quad a_4(d)=n^{8\beta}nd^{-3/2}.
\end{equation}
For $r_2(d)<1$ let $A_2(d)=\{|\rho_2|>r_2(d)\}$.  For $r_4(d)<1$ let
\[
 A_4(d)=\left\{\max_{a<b}|\rho_2(X^a,X^b)|\le r_2(d),\quad
                    |\rho_4(X^1,X^2,X^3,X^4)|>r_4(d)\right\}.
\]
The displayed four-replica probability below is an unconditional
$\mu^{\otimes4}$ probability.

\begin{lemma}[Static support]\label{hist:static-support}
There is a smooth function $\chi_n(\theta)\in[0,1]$ with the following
properties.  On its support, and on the support of any of its derivatives
of order at most $K$, the matrix satisfies
\begin{equation}\label{hist:static-conditions}
 \begin{gathered}
 \mu^{\otimes2}(A_2(d))\le e^{-a_2(d)},\qquad
 \mu^{\otimes4}(A_4(d))\le e^{-a_4(d)},\qquad d\in\mathscr D_n,\\
 n^{-6}\sum_{i<j}\theta_{ij}^2\le2.
 \end{gathered}
\end{equation}
An assertion involving an omitted event is void.  For some $b=b(\beta)>0$,
every derivative of $\chi_n$ of positive order at most $K$ in one coefficient
$\theta_e$, and every differentiated factor in its product construction,
is $O(n^{-b})$, uniformly in $e$ and $\theta$.  For Gaussian couplings,
$\chi_n=1$ with probability tending to one.
\end{lemma}

\begin{proof}
Choose a fixed smooth increasing function $f$ that is zero on $(-\infty,1]$
and one on $[2,\infty)$.  For each nonempty $A_j(d)$ use the factor
$f(-\log\mu^{\otimes j}(A_j(d))/a_j(d))$, and omit empty events.  Multiply
these factors and a smooth function of $n^{-6}\sum_e\theta_e^2$ that is one
below $1$ and zero above $2$.  This gives \eqref{hist:static-conditions} and
a plateau equal to one when each probability there is at most $e^{-2a_j(d)}$
and the size statistic is at most one.

To check derivatives, the $k$th derivative of the logarithm of a nonempty
replica probability in $\theta_e$ is the difference of two cumulants of
$\sum_{a=1}^jX_i^aX_j^a$, under the restricted and unrestricted replica
probabilities.  The variable is bounded by $j\le4$, so every fixed-order
cumulant is bounded by a constant depending only on its order.  Since
$a_2(d)\ge c_0 n^{3\beta}$ and $a_4(d)\ge c_0 n^{8\beta}$ on
$\mathscr D_n$, every derivative of positive order of a replica factor gains at
least $n^{-3\beta}$.  Derivatives of the size factor gain at least $n^{-3}$:
where a derivative is nonzero, $|\theta_e|\le\sqrt2 n^3$.  There are
$O(\log n)$ factors.  The product rule now proves the stated bound, for
example with any sufficiently small $b<\min(3\beta,3)$.

We prove the Gaussian assertion with the diagonal and eigenvector sign
randomization fixed as in Section~\ref{sec:edge-estimates}.  On spectral
events of limiting probability as close to one as desired,
Lemma~\ref{edge:spherical-tails} gives, for the uncapped zero-field spherical
law and a sufficiently small fixed $\rho_0>0$,
\[
 \Prob_{\rm sph}^{\otimes2}\{ |\rho_2|>w\}
       \le C e^{-c_0nw^3},\qquad C_0p\le w\le\rho_0,
\]
where one can take any permitted reference scale with $np^3$ above the
tight threshold.  The smallest pair threshold $r_2(d)$ in
\eqref{hist:static-scales} is $\gg n^{-1/3}$, so such a $p$ is available
below all the nonvacuous shrinking pair thresholds.  In a band
$z<|\rho_2|\le2z$, truncated at $\rho_0$, we apply the weighted
two-replica overlap-cell consequence of Lemma~\ref{edge:gram-cell}.
On the retained event $\|W\|\le3$, the logarithm of the positive zero-field
angular weight has overlap derivative $O(n)$, uniformly over the
orthonormal two-frame.  It therefore oscillates by at most a fixed constant
on each cell of width $2/n$.  Summing the weighted cell inequality over the
band enlarges its spherical endpoints by at most $1/n$ and costs at most
$\exp(Cnz^4+C\log n)$ after increasing $C$.  For the pair thresholds $w$
at issue, $z\ge w\gg n^{-1/3}$, so the enlargement is absorbed by
decreasing the tail constant $c_0$.  Shrink $\rho_0$ so that this
cost is bounded by $\exp(c_0nz^3/2+C\log n)$.  Summing the dyadic bands from
$w$ to $\rho_0$ therefore bounds the Haar mean of the restricted cube
partition function, divided by the squared spherical partition function,
by $\exp(-c_1nw^3)$ for all the thresholds in question.  Here
$nw^3\ge c n^{6\beta}$ absorbs $O(\log n)$.  The likelihood ratio
$L_\circ=Z_{\rm cube}/Z_{\rm sp}$ is at least $1/2$ with probability
tending to one on these restrictions, by Equation~\eqref{edge:haar-bias}.  Markov's
inequality consequently proves the pair plateau, because
\[
 \frac{nr_2(d)^3}{a_2(d)}=n^{3\beta}\longrightarrow\infty.
\]
For overlaps bounded away from zero, the fixed-overlap estimate in
\cite[Lemma~8.2]{FUA} gives an $e^{-c n}$ tail with Gaussian probability tending to
one.  This pays $2a_2(d)=2n^{1-3\beta}r_2(d)^3=o(n)$ for any such
nonvacuous threshold.  %

Here is the corresponding four-replica calculation.  If $r_4(d)<1$, then
$r_2(d)\le n^{-14\beta/3}=o(1)$ and
$r_2(d)^2/r_4(d)\le n^{-6\beta}$.  For one row of four independent uniform
signs, let $Y\in\{-1,1\}^6$ be its six pair products and let $H$ be its
four-product.  For $n$ independent rows write
$\bar Y=n^{-1}\sum_iY_i$ and $\bar H=n^{-1}\sum_iH_i$.  Put
$\Lambda(v)=\log\E e^{v\cdot Y}$, and let $\Prob_v$ be the row probability
with density $e^{v\cdot Y-\Lambda(v)}$, with expectation $\E_v$.  For
several rows the same notation denotes its product probability.  Around the origin,
\[
 \Lambda(v)=\tfrac12|v|^2+O(|v|^3),\qquad
 \E_v H=O(|v|^2).
\]
Indeed the six coordinates of $Y$ are centered and have identity covariance, while
$\E H=\E(HY)=0$; Taylor's theorem on this finite space gives both uniform
remainders.  If the empirical pair vector is the lattice value $u$ with
$|u|\le\sqrt6 r_2(d)$, put $E_d=\{|\bar H|>r_4(d)\}$.  The change of
measure gives
\[
 \Prob_{\rm cube}^{\otimes4}\{\bar Y=u,\ E_d\}
 =e^{-n|u|^2+n\Lambda(u)}
       \Prob_u\{\bar Y=u,\ E_d\}
 \le e^{-n|u|^2+n\Lambda(u)}\Prob_u(E_d).
\]
Under $\Prob_u$ the rows remain independent.  Since
$|\E_uH|\le r_4(d)/2$ for large $n$, Hoeffding's inequality
\cite{Hoeffding1963} bounds the last probability by
$2e^{-nr_4(d)^2/8}$.  Together with the expansion of $\Lambda$, this gives
\begin{equation}\label{hist:cube-quartet}
 \Prob_{\rm cube}^{\otimes4}\{\bar Y=u,\ |\bar H|>r_4(d)\}
 \le 2\exp\{-\tfrac n2|u|^2+Cn|u|^3-cn r_4(d)^2\}.
\end{equation}

For completeness, the reverse Gram cost is also explicit.  Put
$G_{aa}=1$ and $G_{ab}=u_{ab}$.  The Gram density of four independent
uniform points of the sphere of radius $\sqrt n$ is a polynomial
normalization times $\det(G)^{(n-5)/2}$ on positive correlation matrices.
For $|u|=o(1)$,
$\log\det G=-|u|^2+O(|u|^3)$.  A cell of width $c/n$ about $G$ thus has
probability at least
\[
 \exp\{-\tfrac n2|u|^2-Cn|u|^3-C\log n\}.
\]
At a fixed spectrum with bounded norm, let $F\in\R^{n\times4}$ be a uniform
orthonormal frame and put
\[
 \mathcal A_\lambda(G)=Z_{\rm sp}^{-4}\E_F
      \exp\{(n/2)\Tr(F^{\mathsf T}\diag(\lambda)F G)\}.
\]
The exponent has Lipschitz constant $O(n)$ in $G$, so the integrand changes
by at most a factor $e^{O(1)}$ across the cell.  The integral of
$\mathcal A_\lambda$ against the spherical Gram probability equals one.
Restricting it to the cell therefore bounds the Haar mean at Gram $G$ by
\[
 \exp\{\tfrac n2|u|^2+Cn|u|^3+C\log n\}.
\]
Multiply this by \eqref{hist:cube-quartet} and sum the at most $(n+1)^6$
possible Grams.  The remaining error $Cnr_2(d)^3+C\log n$ is negligible
relative to $nr_4(d)^2$, since
\[
 nr_2(d)^3=n^{6\beta}nd^{-3/2},\qquad
 nr_4(d)^2=n^{20\beta}nd^{-3/2}.
\]
On $L_\circ\ge1/2$, division by the cube partition function changes the
four-replica numerator by at most $16$.  Markov's inequality proves its
plateau $e^{-2a_4(d)}$, since $nr_4(d)^2/a_4(d)=n^{12\beta}$.  A union over
$O(\log n)$ scales is harmless.  Finally, the Gaussian mean of
$n^{-6}\sum_e\theta_e^2$ is $O(n^{-5})$.  The size plateau follows from
Markov's inequality.  Enlarging the spectral restrictions completes the
proof that $\chi_n=1$ with probability tending to one.
\end{proof}

Only this last verification used Gaussian frame averages.  In the remainder
of this section and in the score estimates, $\theta$ is any deterministic
matrix satisfying \eqref{hist:static-conditions}.

\subsection{The tree kernels and their survival probabilities}

For each leg $I=[a,b]$, repeatedly bisect time intervals until their lengths
belong to $[1,2]$.  (Changing the harmless endpoint convention when a length
is exactly $2$ leaves all statements below unchanged.)  A node and its two
children are denoted by $I$ and $I^-,I^+$, in chronological order.  Define
\begin{equation}\label{hist:node-level}
 N_I=N(|I|),\qquad N(d)=n^\beta(T/d)^{3/2}=n^{1+\beta}d^{-3/2}.
\end{equation}
All these levels lie between $c_0n^\beta$ and $C_0n^{1+\beta}$.

\begin{lemma}[A smooth cap]\label{hist:cutoff-hazards}
For every level in \eqref{hist:node-level} there is a smooth
$W_N:\R\to(0,1)$ such that
\begin{equation}\label{hist:cap-and-plateau}
 e^lW_N(l)\le e^{CN},\qquad
 1-W_N(l)\le e^{-(\log n)^2/2}\quad(l\le N/2).
\end{equation}
For $1\le k\le K$ its derivative in $l$ satisfies globally
\begin{equation}\label{hist:global-ratios}
 \frac{|W_N^{(k)}|}{W_N}+\frac{|W_N^{(k)}|}{1-W_N}\le C_K,
 \qquad
 \frac{|W_N^{(k)}|}{1-W_N}\le\frac{C_K(\log n)^{2K}}N.
\end{equation}
For each fixed $L$, wherever $W_N(l)\ge n^{-L}$ one also has
\begin{equation}\label{hist:survival-hazard}
 \frac{|W_N^{(k)}(l)|}{W_N(l)}
 \le \frac{C_{K,L}(\log n)^{2K}}N\bigl(1-W_N(l)\bigr).
\end{equation}
Moreover, uniformly over these levels and for sufficiently large $n$,
\begin{equation}\label{hist:log-cap-growth}
 |\log W_N(l)|+|\log(1-W_N(l))|
       \le C\bigl(1+N+|l|+(\log n)^2\bigr).
\end{equation}
\end{lemma}

\begin{proof}
Let $\psi$ be smooth and nondecreasing, zero on $(-\infty,2]$, with
$0\le\psi'\le1$ and $\psi'=1$ on $[3,\infty)$.  Put
\[
 a_N(z)=(\log n)^2(z-1)+3N\psi(z),\qquad
 W_N(l)=\bigl(1+e^{a_N(l/N)}\bigr)^{-1}.
\]
For $l\le3N$ the first inequality in \eqref{hist:cap-and-plateau} is
immediate.  For $z=l/N\ge3$, $\psi(z)\ge z-3$, so
$l-a_N(z)\le3N$; the same inequality follows from $W_N\le e^{-a_N}$.
If $l\le N/2$, then $a_N(l/N)\le-(\log n)^2/2$, proving the plateau.

All $l$-derivatives of $a_N(l/N)$ of positive order at most $K$ are bounded:
the first is $(\log n)^2/N+3\psi'(l/N)$ and the others are
$3N^{1-k}\psi^{(k)}(l/N)$.  Derivatives of the logistic function have a
factor $W_N(1-W_N)$ times a bounded polynomial in $W_N$, which proves the
first part of \eqref{hist:global-ratios}.  If $W_N\ge n^{-L}$, then
$a_N\le L\log n+O(1)<(\log n)^2$; hence $l/N<2$ for large $n$.  In this
region $a_N$ is affine, and its slope in $l$ is $(\log n)^2/N$.
The same derivative formula proves \eqref{hist:survival-hazard}.  It also
proves the second part of \eqref{hist:global-ratios} for $l/N\le2$.  In the
remaining region $W_N\le e^{-(\log n)^2}$, so the first global bound with
its logistic factor gives the second bound there as well, since $N$ is
polynomial in $n$.  Finally the two logarithms are bounded by
$C(1+|a_N(l/N)|)$.  Since $\psi$ is Lipschitz and
$(\log n)^2/N\le1$ for large $n$, this proves \eqref{hist:log-cap-growth}.
\end{proof}

Let $p_d(x,y)=P_d(x,y)/\mu(y)$ be the true transition density.  Define
the pre-density $\bar k_I$, the survival factor $w_I$, and the post-density
$k_I$, all relative to the normalized probability $\mu$, by
\begin{equation}\label{hist:kernel-recursion}
 \begin{aligned}
 \bar k_I(x,y)&=p_{|I|}(x,y) &&\text{if $I$ is a leaf},\\
 \bar k_I(x,y)&=\sum_z\mu(z)k_{I^-}(x,z)k_{I^+}(z,y)
                                      &&\text{otherwise},\\
 w_I(x,y)&=W_{N_I}(\log\bar k_I(x,y)),\qquad
 k_I(x,y)=\bar k_I(x,y)w_I(x,y).
 \end{aligned}
\end{equation}
All these numbers are positive.  In particular, the logarithms are defined.
The transition $k_I(x,y)\mu(y)$ is a subprobability.  At a fixed leg we
may equivalently sample an ordinary marked heat-bath path and retain it with
the product of its $w_I(X_a,X_b)$ over all nodes.  Apply checks at their
right endpoints, ordering equal right endpoints from the smallest interval
to the largest.  The equivalence follows inductively from the Markov
property and \eqref{hist:kernel-recursion}.

\begin{lemma}[Deterministic kernel consequences]\label{hist:kernel-caps}
For every node $I$ of length $d$,
\begin{equation}\label{hist:kernel-bounds}
 0<k_I(x,y)\le\bar k_I(x,y)\le p_d(x,y),\qquad k_I(x,y)\le e^{CN(d)}.
\end{equation}
Both $k_I$ and $\bar k_I$ are symmetric in $(x,y)$.  If an interval of
length $h$ lies inside a checked tree, with $h$ larger than a fixed constant,
the weighted transition across it has density at most $e^{C'N(h)}$ after
all weights other than one contained complete node are discarded.  The same
assertion holds when the top check of the tree is omitted or is a failure
factor $1-w_I$.
\end{lemma}

\begin{proof}
Positivity and domination by $p_d$ follow by induction and the semigroup
property of $P_d$.  The cap is \eqref{hist:cap-and-plateau}.  At a split the
two children have equal lengths and use the same rule, so their kernels are
identical.  The convolution of this symmetric kernel with itself is
symmetric; multiplication by the endpoint function $W_N(\log\bar k_I)$
preserves symmetry.  This starts from detailed balance of $p_d$.

Every subinterval of length $h$ larger than a fixed constant contains a
tree node $J$ of length at least $c_1h$.  One may require $J$ to be a proper
descendant at the cost of reducing $c_1$.  Keep the checks in $J$ and discard
the other weights, all of which lie in $[0,1]$.  The remaining transition
is an ordinary kernel, followed by $k_J$, followed by an ordinary kernel.
Since ordinary kernels preserve $\mu$ and constants, its density is at most
$e^{CN(|J|)}\le e^{C'N(h)}$.  A discarded failure or top factor also lies
in $[0,1]$.
\end{proof}

\begin{lemma}[Gaussian survival]\label{hist:gaussian-survival}
For every $A<\infty$ and $\varepsilon>0$, there are Gaussian disorder
events of limiting lower probability at least $1-\varepsilon$ on which
$\chi_n=1$ and the mass lost by the checks in any complete query is at most
$n^{-A}$.  This bound is uniform over its query times in the fixed compact
set and over its initial probability on the cube.
\end{lemma}

\begin{proof}
The log-moment assertion of Theorem~\ref{thm:warming} says that, for each fixed
integer $p\ge1$ and sufficiently small fixed $\xi>0$, on disorder restrictions
of arbitrarily high limiting probability,
\[
 \sup_x\left(1+\sum_yP_d(x,y)\log_+^p p_d(x,y)\right)^{1/p}
 \le C_{p,\xi}\left(1+(T/d)^{1/(2/3-\xi)}\right),\quad 1\le d\le CT.
\]
Choose $\xi$ so that $\xi/(2/3-\xi)<\beta/2$.  Division by $N(d)$ shows
that the right side is at most $C_{p,\xi}n^{-\beta/2}N(d)$ uniformly in
this range: the largest possible extra factor is
$T^{1/(2/3-\xi)-3/2}=n^{\xi/(2/3-\xi)}$.  Markov's inequality therefore
gives
\[
 \sup_xP_d(x,\{y:\log p_d(x,y)>N(d)/2\})
       \le C_{p,\xi}n^{-p\beta/2}.
\]
By \eqref{hist:kernel-bounds}, $\bar k_I\le p_{|I|}$.  At each check, its
failure probability under an ordinary path is thus bounded by the last
display plus $e^{-(\log n)^2/2}$.  The probability of a first failure is
at most the sum of these probabilities, because all preceding weights are
at most one.  There are $O(T)$ nodes.  Choose $p$ large after $A$ and use
Lemma~\ref{hist:static-support} for the plateau.  All estimates were uniform
in the starting state and in $d$, which proves the asserted uniformity.
\end{proof}

\subsection{Tested histories and their conditional interiors}

We record precisely which changes of path measure will be used.  An
ordinary marked path includes the time, site, and new sign at every ring,
including a resampling that keeps the current sign.  Its law on a time
interval $I$ from $x$ is denoted by $\mathsf P_x^I$.

\begin{definition}[Tested path law]\label{hist:tested-law}
A checked history is one of the following subprobabilities, with initial
probability either $\mu$ or the uniform cube probability:
\begin{enumerate}
\item[(i)] an entire query path with every survival factor in its legs;
\item[(ii)] the prefix through a specified first failing check $F$, with all
preceding survival factors and the factor $1-w_F$ at that check.
\end{enumerate}
The check order is the one following \eqref{hist:kernel-recursion}; the
choice of $F$ and all tree shapes is fixed independently of the path and
of $e$.  In case (ii), the maximal checked intervals of the prefix form a
forest: the last is $F$, whose descendants survive, and every preceding
root and its descendants survive.  Designate the initial state, all leg
boundaries, and all boundaries of these maximal roots as its external
endpoints, retaining only boundaries present in a stopped prefix.  Multiply
the history by a nonnegative function $\Phi$ of these
endpoints with $\|\Phi\|_\infty\le n^{M_+}$, where $M_+$ is fixed and
$\Phi$ is the same for every $e\in\mathscr E_n$.  When the resulting mass
$m_\Phi$ is at least $n^{-M_-}$ for a fixed $M_-$, its normalization is a
tested path law, denoted by $\mathsf Q_\Phi$.

Products of a fixed number of such histories use independent ordinary
paths and a separate external test on each history.  Signed tests are
decomposed into positive and negative parts separately.  A statement that
holds with arbitrarily high inverse-polynomial probability means that for
every fixed $A$ its exception has $\mathsf Q_\Phi$ probability at most
$n^{-A}$ for all sufficiently large $n$, uniformly in the fixed choices
above and in disorder satisfying \eqref{hist:static-conditions}.
\end{definition}

There are at most polynomially many nodes.  In the last, possibly partial,
leg of a first-failure history, its preceding maximal roots are the earlier
siblings encountered in the binary subdivision leading to $F$.  In their
chronological order their lengths are nonincreasing and the last such
length is at least $|F|$.  Complete earlier legs have lengths comparable
to $T$.  These facts also show that earlier maximal roots have lengths at
least a fixed multiple of later ones, with the multiple depending only on
the query set.  The density of $\mathsf Q_\Phi$ relative to the ordinary
forward law with the same initial probability is at most
$n^{M_++M_-}$; future endpoints are not revealed in that forward law's
natural filtration.

For a node $I=[a,b]$, write
$V_I=\prod_{J\subsetneq I}w_J(X_{\partial J})$, with the product over its
proper descendants.  The normalized interior law with boundary $(x,y)$ is
\begin{equation}\label{hist:bridge-definition}
 \mathsf C_I^{x,y}(d\omega)
 =\frac{V_I(\omega)\ind_{\{X_b=y\}}}{\mu(y)\bar k_I(x,y)}
             \mathsf P_x^I(d\omega).
\end{equation}
The denominator equals the numerator's integral by the kernel recursion.
This definition excludes the top check.  A top survival, failure, or omitted
check is constant after its boundary is fixed.

\begin{lemma}[Conditional interiors and prefix bounds]\label{hist:conditional-interior}
Let $I$ be a fixed node of a checked root in a tested path law.  Given the marked
history outside the interior of $I$ and its two boundary states, the
conditional law inside $I$ is $\mathsf C_I^{X_a,X_b}$.  In particular,
resampling this interior independently from that law preserves the tested
marginal.  For an internal node its midpoint $z$ under $\mathsf C_I^{x,y}$
has probability
\begin{equation}\label{hist:midpoint-probability}
 \pi_I^{x,y}(z)=
   \frac{\mu(z)k_{I^-}(x,z)k_{I^+}(z,y)}{\bar k_I(x,y)},
\end{equation}
and conditional on $z$ the two interiors are independent with laws
$\mathsf C_{I^-}^{x,z}$ and $\mathsf C_{I^+}^{z,y}$.

For every $A$, there is $L=L(A,M_+,M_-)<\infty$ such that
$\bar k_I(X_a,X_b)\ge n^{-L}$ outside an endpoint event of probability
$n^{-A}$ under any such tested law.  If $d=|I|$ exceeds a fixed constant,
then for these boundaries the density of its first two thirds under
$\mathsf C_I^{x,y}$ relative to the ordinary unconditioned prefix from $x$
is at most
\begin{equation}\label{hist:prefix-density}
 n^L e^{C'N(d)}.
\end{equation}
The reversed last two thirds have the same bound relative to the ordinary
prefix from $y$ in reversed time.
\end{lemma}

\begin{proof}
The tree is laminar: every check whose interior intersects the interior of
$I$ is a descendant or an ancestor of $I$.  Disjoint siblings may share an
endpoint, but their factors are fixed by the exterior and boundary data.
The ancestor factors, the top factor, and the external test are fixed by
the same data.  The remaining factors are exactly
$V_I$.  The Markov property gives \eqref{hist:bridge-definition} as the
conditional law, and then gives \eqref{hist:midpoint-probability} and the
conditional independence.  This proves the resampling assertion as well.
It applies to tree nodes; an arbitrary subinterval can have crossing checks.

At a fixed start $x$,
\[
 \sum_y\mu(y)\bar k_I(x,y)\ind_{\{\bar k_I(x,y)<\delta\}}\le\delta.
\]
Integrate the node first in the tested history and discard all other weights.
The bound on $\Phi$ and the lower bound on its mass show
$\mathsf Q_\Phi\{\bar k_I(X_a,X_b)<\delta\}\le n^{M_++M_-}\delta$.
This proves the endpoint assertion by choosing $\delta=n^{-L}$.

Let $\mathcal F_t^\to$ contain just the marked prefix of the node from $a$
to $t$.  For $t-a\le2d/3$, its conditional density is exactly
\[
 \frac{d(\mathsf C_I^{x,y}|_{\mathcal F_t^\to})}
      {d(\mathsf P_x^I|_{\mathcal F_t^\to})}
 =\frac{\mathsf E_x[V_I\ind_{\{X_b=y\}}\mid\mathcal F_t^\to]}
        {\mu(y)\bar k_I(x,y)}.
\]
The interval remaining after $t$ contains a proper node of length at least
$c_1d$.  Keep its checks and discard all others in the numerator.  By
Lemma~\ref{hist:kernel-caps}, that numerator is at most $\mu(y)e^{C'N(d)}$.
This proves \eqref{hist:prefix-density}.  Detailed balance for marked paths
includes self-resampling rings, and the proper-descendant weights are
invariant under reversal because equal-length children use identical
symmetric kernels.  Reversal therefore sends $\mathsf C_I^{x,y}$ to
$\mathsf C_I^{y,x}$ on the reflected node, proving the last assertion.
No reversal of the entire stopped forest is needed.
\end{proof}

\begin{lemma}[Overlaps along a checked history]\label{hist:history-overlaps}
With arbitrarily high inverse-polynomial probability under a tested path
law, simultaneously for times in any one of its checked root intervals,
\begin{equation}\label{hist:path-overlaps}
 |\rho_2(X(t_1),X(t_2))|\lesssim_\beta(1+m)^{-1/2},\qquad
 |\rho_4(X(t_1),X(t_2),X(t_3),X(t_4))|
                    \lesssim_\beta(1+m)^{-3/4}.
\end{equation}
In each assertion the times are increasingly ordered, and $m$ is the
minimum of their consecutive gaps and of the gap from the root's start to
$t_1$.  The assertion includes a failing root.  It also holds for the
full history obtained by each of a fixed finite collection of interior
resamplings in Lemma~\ref{hist:conditional-interior}, separately within
each such history.  All times in an overlap are taken from that one history.
\end{lemma}

\begin{proof}
First use deterministic times with $m$ larger than a fixed constant.
Inside each gap choose a contained proper checked node of length comparable
to $m$.  They are disjoint.  Dropping all other weights and applying
Lemma~\ref{hist:kernel-caps} shows that the joint law of the observed states,
given any initial state, is bounded by their product $\mu$ law times
$\exp(C_2N(m))$.  This estimate counts the retained weights without
conditioning on survival.  Apply \eqref{hist:static-conditions} at a
dyadic scale comparable to $m$, first to all relevant pairs, and then to
the quartet on the event that its pairs obey the same threshold.  The
ratios $a_2(m)/N(m)=n^{2\beta}$ and $a_4(m)/N(m)=n^{7\beta}$ tend to
infinity.  Their exponents are at least a positive power of $n$, so they
pay the cap and any polynomial test density with any prescribed
inverse-polynomial precision.  A vacuous threshold or bounded $m$ gives
the bound directly from $|\rho_j|\le1$.

Use a time mesh of sufficiently small inverse-polynomial spacing.  The
probability that any mesh cell has two rings is arbitrarily small after
paying the polynomial test density: under ordinary paths the total clock
rate is $n$, and the union bound is $O(n^2T)$ times the spacing.  A single
ring changes either overlap by at most $2/n$ per affected time.  This is
smaller than the displayed thresholds, whose smallest quartet value has
order $n^{-1/2+10\beta}$.  A union over the polynomially many grid tuples
proves the simultaneous assertion.  A resampled interior has the same
tested marginal by Lemma~\ref{hist:conditional-interior}; a union proves the
last claim for any fixed number of copies.
\end{proof}

\subsection{Local differentiability and normalizations}

We finish by recording a coarse derivative bound, used to control exceptional
histories.  For $e=\{i,j\}$ write $q_e(x)=x_ix_j$ and perturb the interaction
$\theta_e$ by $\vartheta$, putting $u=\tanh\vartheta$.  The new stationary
probability is exactly
\begin{equation}\label{hist:stationary-perturbation}
 \mu_u(x)=\mu(x)\frac{1+u q_e(x)}{1+u\mu q_e},\qquad |u|<1.
\end{equation}
Let $t_i(x)=\mu[X_i\mid X_{-i}=x_{-i}]$.  At a ring of site $i$ with new
sign $\zeta$, the likelihood ratio of the perturbed mark is
\begin{equation}\label{hist:ring-likelihood}
 M_{e,i,\zeta}(u)=\frac{1+u\zeta x_j}{1+u x_jt_i(x)};
\end{equation}
interchange $i,j$ at the other site, and use one at all remaining sites.
The site rates stay equal to one.  Thus the full path likelihood
$L_{I,e}(u)$ is the product of \eqref{hist:ring-likelihood} at its rings.
The formula includes a mark that retains the current sign.

Superscript $u$ on $\bar k_I,k_I,w_I$ denotes the construction
\eqref{hist:kernel-recursion} for the perturbed dynamics and its normalized
stationary probability $\mu_u$.  Throughout, $[u^l]$ means the $l$th Taylor
coefficient at zero (the derivative divided by $l!$).  For the pre-density,
the exact change of measure is
\begin{equation}\label{hist:pre-ratio}
 \frac{\bar k_I^u(x,y)}{\bar k_I(x,y)}
 = (1+u\mu q_e)\,
 \mathsf C_I^{x,y}\left[
    \frac{L_{I,e}(u)}{1+u q_e(y)}
      \prod_{J\subsetneq I}\frac{w_J^u(X_{\partial J})}{w_J(X_{\partial J})}
                         \right].
\end{equation}
The product is over all proper descendants; it contains only survival
factors.  In particular $1+u\mu q_e$ is required by normalization of the
perturbed stationary probability.

\begin{lemma}[Coarse coefficients and moments]\label{hist:crude-ratios}
On the size support in \eqref{hist:static-conditions}, for all nodes,
boundaries and $e$, and $1\le l\le K$,
\[
 \left|[u^l]\frac{\bar k_I^u}{\bar k_I}\right|
 +\left|[u^l]\frac{k_I^u}{k_I}\right|
 +\left|[u^l]\frac{w_I^u}{w_I}\right|
 +\left|[u^l]\frac{1-w_I^u}{1-w_I}\right|\le n^{C_K}.
\]
The exponent $C_K$ is independent of $\beta$ in a fixed sufficiently small
range and of any later number of conditional copies.  Absolute fixed-order
moments of clock counts and raw path coefficients are polynomial under
tested laws and their resampled marginals.  Conditional on an arbitrary
boundary pair of one node, the same moments under $\mathsf C_I^{x,y}$ are
bounded by a polynomial (with exponent allowed to depend on the moment
order).  Finally, conditional on arbitrary designated external endpoints,
the absolute normalized coefficient through total order $K$ of an entire
checked or first-failure history is bounded by $n^{C_K}$, after enlarging
$C_K$.  Thus its unnormalized coefficient with any nonnegative external
endpoint test held fixed under the interaction perturbation is bounded by
$n^{C_K}$ times the test's base mass, however
small that mass is.
\end{lemma}

\begin{proof}
The size support gives $\sum_e|\theta_e|\le Cn^4$, hence
$\min_x\mu(x)\ge e^{-Cn^4}$.  By Cauchy--Schwarz, each absolute one-site
field is at most $Cn^{7/2}$.
On a leaf, prescribe one ring of each site in successive subintervals of
length $1/n$, with its sign set to the desired endpoint, and no other rings.
The clock event has probability $e^{-n|I|}n^{-n}$, and each prescribed mark
has probability at least $\tfrac12e^{-Cn^{7/2}}$.  Thus every leaf transition
probability is at least $e^{-Cn^5}$, uniformly in its endpoints.  Its
logarithmic density has absolute value at most $Cn^5$.  The last assertion of
Lemma~\ref{hist:cutoff-hazards} and the kernel recursion propagate this property
up the tree: composing two positive kernels costs a fixed multiple of the
previous logarithmic bound plus a polynomial, and applying $W_N$ has the
same property.  The depth is $O(\log n)$, so all these absolute logarithms
remain polynomial, with a fixed exponent.  In particular
$\mu(y)\bar k_I(x,y)\ge e^{-n^{C_0}}$ for a fixed $C_0$.

The total number of ordinary rings on a node is Poisson with parameter
$n|I|\le CnT$.  Its tail above any sufficiently large polynomial is smaller
than $e^{-2n^{C_0}}$ times any inverse power of $n$.  Dividing by the last
lower bound and using $V_I\le1$ proves polynomial conditional count moments
for every boundary.  Each fixed Taylor coefficient of the ring product is
a polynomial of fixed degree in its number of rings, since each single-ring
coefficient in \eqref{hist:ring-likelihood} is bounded.  This proves the
claimed raw moments, both conditionally and under the polynomially tilted
tested laws.  The copy statement follows from their preserved marginals.

At a leaf, \eqref{hist:pre-ratio} has an empty product; the preceding count
bound and the bounded coefficients of $(1+u\mu q_e)/(1+u q_e(y))$ give
the coarse bound for its pre-density.  Suppose normalized coefficients of
order $l\le K$ at one tree level are bounded by $D^l$.  For a composition,
division by its base value writes its ratio as a conditional average of the
product of the two child ratios and $\mu_u(z)/\mu(z)$.  Its order-$l$
coefficient is therefore at most $(C_K(D+1))^l$.  Taylor's chain rule for
$W_N(\log\bar k_I^u)$ uses the global ratios
\eqref{hist:global-ratios}; coefficients of the logarithm are polynomials
in nonconstant pre-ratio coefficients, each product having total order $l$.
The same bound applies to either survival or failure ratios and to the
post-density.  Starting with a polynomial $D$ at the leaves and iterating
$D\mapsto C_K(D+1)$ for $O(\log n)$ levels proves these coefficient bounds.  Formula
\eqref{hist:pre-ratio} itself follows by inserting the path likelihood and
then dividing by $\mu_u(y)/\mu(y)$ in the definition of the pre-density.

To check the last forest assertion, fix its designated external endpoints.
The interiors of its maximal roots are the corresponding conditional laws
\eqref{hist:bridge-definition}, and their top survival or failure factors
are fixed.  Expand a coefficient of total order at most $K$ by the product
rule.  A differentiated node contributes one of the normalized ratios just
bounded; the raw ring factors have the conditional polynomial moments
proved above.  There are polynomially many nodes and at most $K$ positive
orders, so the number and magnitude of the resulting terms are bounded by
$n^{C_K}$.  Stationary initialization, if present, contributes only the
bounded coefficients in \eqref{hist:stationary-perturbation}.  Integrating
this conditional bound against the nonnegative endpoint test and its base
endpoint measure proves the asserted factor of the base mass.
\end{proof}

Here and below, estimates on a set of arbitrarily small inverse-polynomial
probability can be used inside conditional averages.  The precise elementary
rule is as follows.  If $\mathsf Q(E^c)\le n^{-A'}$ and
$\|Y\|_{L^p(\mathsf Q)}\le n^C$ for some $p>1$, then for any boundary
$\sigma$-field $\mathcal B$,
\[
 \mathsf Q\left\{\mathsf E_{\mathsf Q}
       [|Y|\ind_{E^c}\mid\mathcal B]>n^{-A}\right\}
 \le n^{A+C-(1-1/p)A'}.
\]
Choose $A'$ after $A$ to make this as small as required.  The coarse bounds
of Lemma~\ref{hist:crude-ratios} supply these moments for all coefficients used
below; a fixed number of copies is handled by H\"older's inequality.
For a tested mass below a chosen inverse-polynomial precision, use instead
the final forest assertion of Lemma~\ref{hist:crude-ratios} before
normalization: its unnormalized coefficient is at most $n^{C_K}$ times
that mass.  Polynomially many deterministic nodes,
grid points, or choices of any fixed number of copies are handled by a union
bound.  These conventions never give a uniform assertion for every boundary
pair when the sharper estimate was proved only for typical boundaries.
